\documentclass[12pt]{article}

\usepackage{amsthm}
\usepackage{amsmath}
\usepackage{amsfonts}
\usepackage{mathrsfs}
\usepackage{amssymb}
\usepackage{units}
\usepackage{graphicx}
\usepackage{tikz-cd}
\usepackage{nicefrac}
\usepackage{hyperref}
\usepackage{bbm}
\usepackage{color}
\usepackage{tensor}
\usepackage{tipa}
\usepackage{bussproofs}
\usepackage{ stmaryrd }
\usepackage{ textcomp }
\usepackage{leftidx}
\usepackage{afterpage}
\usepackage{varwidth}
\usepackage{physics}

\newcommand\blankpage{
    \null
    \thispagestyle{empty}
    \addtocounter{page}{-1}
    \newpage
    }

\graphicspath{ {images/} }

\theoremstyle{plain}
\newtheorem{thm}{Theorem}[subsection] % reset theorem numbering for each chapter
\newtheorem{proposition}[thm]{Proposition}
\newtheorem{lemma}[thm]{Lemma}
\newtheorem{fact}[thm]{Fact}
\newtheorem{cor}[thm]{Corollary}
\newtheorem{post}[thm]{Postulate}

\theoremstyle{definition}
\newtheorem{defn}[thm]{Definition} % definition numbers are dependent on theorem numbers
\newtheorem{exmp}[thm]{Example} % same for example numbers
\newtheorem{notation}[thm]{Notation}
\newtheorem{remark}[thm]{Remark}
\newtheorem{condition}[thm]{Condition}
\newtheorem{question}[thm]{Question}
\newtheorem{construction}[thm]{Construction}
\newtheorem{exercise}[thm]{Exercise}
\newtheorem{example}[thm]{Example}
\newtheorem{observation}[thm]{Observation}
\newtheorem{algorithm}[thm]{Algorithm}

\newcommand{\bb}[1]{\mathbb{#1}}
\newcommand{\scr}[1]{\mathscr{#1}}
\newcommand{\call}[1]{\mathcal{#1}}
\newcommand{\psheaf}{\text{\underline{Set}}^{\scr{C}^{\text{op}}}}
\newcommand{\und}[1]{\underline{\hspace{#1 cm}}}
\newcommand{\adj}[1]{\text{\textopencorner}{#1}\text{\textcorner}}
\newcommand{\comment}[1]{}
\newcommand{\lto}{\longrightarrow}

\usepackage[margin=1cm]{geometry}

\title{Quantum computing}
\author{Will Troiani}
\date{August 2020}

\begin{document}

\maketitle
\tableofcontents
\section{Mathematical definitions of physical concepts}\label{Sec:Math_defs}
\begin{defn}
A \textbf{state space} is a complex Hilbert space, a unit vector in $\bb{H}$ is a \textbf{state}. A \textbf{qubit} is the state space $\bb{H} \cong (\bb{C}^2)^{\otimes n}$ for some $n > 0$.\footnote{Standard usage: $\bb{C}^2$ is a single-qubit state space, and $(\bb{C}^2)^{\otimes n}$ is an $n$-qubit state space. See Section~\ref{sec:errata}.} Sometimes to specify that a qubit is $\bb{C}^2$ we will say the qubit is a \textbf{single qubit}.
\end{defn}
The trace of an operator is part of the mathematical formalisation of a \emph{measurement} (Definition \ref{def:measurement}), how the trace becomes part of the formalism hinges crucially on Lemma \ref{lem:trace_evaluation} below.

Throughout, $\bb{H}$ is a fixed state space.
\begin{defn}
The \textbf{trace} of an operator $T: \bb{H} \lto \bb{H}$ is the trace of any (and hence all) matrix representation of $T$.
\end{defn}
The trace of an operator can be written using bra-ket notation: let $\ket{v_1},...,\ket{v_n}$ be an orthogonal basis for $\bb{H}$,\footnote{The formula $\operatorname{Tr}(A)=\sum_i\bra{v_i}A\ket{v_i}$ requires the basis to be \emph{orthonormal}, not merely orthogonal. See Section~\ref{sec:errata}.} and let $A$ be an operator on $\bb{H}$ and write $A\ket{v_j} = a_{1j}\ket{v_1} + \hdots + a_{nj}\ket{v_n}$, then we have
\begin{equation}
\bra{v_i}A\ket{v_j} = a_{ij}
\end{equation}
and so
\begin{equation}
\operatorname{Trace}A = \sum_{i = 1}^n \bra{v_i}A\ket{v_i}
\end{equation}
The trace of $A$ precomposed with a projector has a particularly nice form: let $\ket{\psi}$ be a unit vector in $\bb{H}$ and let $\lbrace \ket{\psi},\ket{v_2},...,\ket{v_n}\rbrace$ be an orthogonal basis for $\bb{H}$. Then
\begin{align*}
\operatorname{Trace}(A\ket{\psi}\bra{\psi}) &= \bra{\psi}A\ket{\psi}\bra{\psi}\ket{\psi} + \sum_{i = 2}^n \bra{v_i}A\ket{\psi}\bra{\psi}\ket{v_i}\\
&= \bra{\psi}A\ket{\psi}
\end{align*}
In a statement:
\begin{lemma}\label{lem:trace_evaluation}
Let $A$ be an operator on a Hilbert space $\bb{H}$ and let $\ket{\psi} \in \bb{H}$ be a unit vector in $\bb{H}$. We have the following formula:
\begin{equation}
\operatorname{Trace}(A\ket{\psi}\bra{\psi}) = \bra{\psi}A\ket{\psi}
\end{equation}
\end{lemma}
\begin{proof}
Use the Gram-Schmidt process to extend $\lbrace \ket{\psi}\rbrace$ to an orthonormal basis, then apply the previous discussion.
\end{proof}
\begin{defn}\label{def:measurement}\footnote{The completeness sum should read $\sum_{m\in\mathcal{M}}M_m^\dagger M_m=I$ (subscripts on each $M_m$). The post-measurement state $M_m\ket{\psi}/\sqrt{p(m)}$ is defined only when $p(m)>0$. See Section~\ref{sec:errata}.}
A \textbf{measurement} on $\bb{H}$ is a finite family of linear operators $\lbrace M_m: \bb{H}\lto \bb{H}\rbrace_{m \in \call{M}}$ satisfying
\begin{equation}
\sum_{m \in \call{M}}M^{\dagger}M = I
\end{equation}
Associated to every measurement and state vector $\ket{\psi}$ there is a value
\begin{equation}
p(m) := \bra{\psi}M^{\dagger}_mM_m\ket{\psi} \in [0,1]
\end{equation}
which is understood as the probability $p(m)$ of outcome $m$ on measurement $M_m$. Notice by Lemma \ref{lem:trace_evaluation} that
\begin{equation}
p(m) = \operatorname{Trace}(M_m^\dagger M_m \ket{\psi}\bra{\psi})
\end{equation}
Also, there is a \textbf{resulting state}, which is the state
\begin{equation}
\frac{M_m\ket{\psi}}{\sqrt{p(m)}}
\end{equation}
understood as the state of $\bb{H}$ after measuring it with $\lbrace M_m\rbrace_{m \in \call{M}}$.
\end{defn}
\begin{defn}\footnote{In quantum mechanics ``projector'' means \emph{orthogonal} projector: $P^2=P$ and $P^\dagger=P$. The bare condition $P^2=P$ defines an algebraic idempotent. See Section~\ref{sec:errata}.}
A linear transformation $P$ is a \textbf{projector} if $P^2 = P$.
\end{defn}
Recall that Hermitian operators $M$ are in particular \emph{Normal}, ie, satisfy $MM^\dagger = M^\dagger M$, and so are diagonalisable (see \cite{analysis}).
\begin{defn}\footnote{The body has ``$m \backslash\text{in}\ \mathcal{M}$''; should read $m\in\mathcal{M}$. Post-measurement state requires $p(m)>0$. See Section~\ref{sec:errata}.}
A \textbf{projective measurement} on $\bb{H}$ is a Hermitian operator $M: \bb{H} \lto \bb{H}$ with spectral decomposition
\begin{equation}
M = \sum_{m \in \call{M}}m P_m
\end{equation}
where $P_m$ is the projector onto the $m$-eigenspace. Associated to value $m \ in \call{M}$ and state $\ket{\psi}$ there is a value
\begin{equation}
p(m) := \bra{\psi} P_m\ket{\psi}
\end{equation}
which is understood as the probability $p(m)$ of outcome $m$ on measurement $P_m$. Also, there is a \textbf{resulting state}:
\begin{equation}
\frac{P_m\ket{\psi}}{\sqrt{p(m)}}
\end{equation}
understood as the state of $\ket{\psi}$ after measuring it with $P_m$.
\end{defn}
\begin{lemma}\footnote{This statement is incorrect: a unitary $U:W\to V$ forces $W=V$ (in finite dimensions), so the conclusion is trivial. The intended hypothesis is that $U$ is an \emph{isometry}. The proof $U'=\operatorname{Proj}_W\circ U$ is also wrong (domain/codomain mismatch). See Section~\ref{sec:errata} for a correct statement and proof.}
Let $W \subseteq V$ be a subspace of a Hilbert space $V$, and let $U: W \lto V$ be a unitary operator. Then $U$ extends to a unitary $U'$ operator on all of $V$.
\end{lemma}
\begin{proof}[Proof sketch]
Define $U' = \operatorname{Proj}_W \circ U$.
\end{proof}
The relationship between projective measurements and measurements is the following:
\begin{proposition}\footnote{The proof switches between $Q\otimes\mathbb{H}$ and $\mathbb{H}\otimes Q$, omits subscripts on $M_m^\dagger M_{m'}$, writes $I_q\otimes\ket{m}\bra{m}$ instead of $I_\mathbb{H}\otimes\ket{m}\bra{m}$, and has a wrong final ket in the inner-product calculation. See Section~\ref{sec:errata} for a clean version.}
Let $\lbrace M_m\rbrace_{m \in \call{M}}$ be a measurement on $\bb{H}$. Then there exists a projective measurement $M$, a state space $Q$, and a unitary operator $U: Q \otimes \bb{H} \lto Q \otimes \bb{H}$ such that the measurement $\lbrace M_m\rbrace_{m \in \call{M}}$ is modelled by $M$ and $U$ (the meaning of \emph{modelled} is provided by the proof).
\end{proposition}
\begin{proof}
Let $Q$ be the free Hilbert space of dimension $|\call{M}|$, the cardinality of $\call{M}$. Let $\ket{\varphi}$ be an arbitrary, fixed vector inside $Q$, and define
\begin{align*}
U: \bb{H} \otimes \operatorname{Span}\ket{\varphi} &\lto \bb{H} \otimes Q\\
\end{align*}
given by 
\begin{equation}
U\ket{\psi}\otimes \ket{\varphi} = \sum_{m \in \call{M}}M_m\ket{\psi}\otimes \ket{m}
\end{equation}
We first prove this is unitary. We calculate:
\begin{align*}
\bra{\psi_1}\otimes \bra{\varphi}U^\dagger U\ket{\psi_2}\otimes\ket{\psi_1} &= \big(U\ket{\psi_1}\otimes\ket{\varphi}\big)^\dagger U\ket{\psi_2}\otimes\ket{\varphi}\\
&= \big(\sum_{m \in \call{M}}M_m\ket{\psi_1}\otimes\ket{m}\big)^\dagger \sum_{m' \in \call{M}}M_{m'}\ket{\psi_2}\otimes\ket{m'}\\
&= \sum_{m \in \call{M}}\bra{\psi}M_m^\dagger \bra{m}\sum_{m' \in \call{M}}M_{m'}\ket{\psi_2}\otimes\ket{m'}\\
&= \sum_{m \in \call{M}}\sum_{m' \in M}\bra{\psi_1}M^\dagger M\ket{\psi_2}\bra{m}\ket{m'}\\
&= \sum_{m \in \call{M}}\bra{\psi_1}M^\dagger M\ket{\psi_2}\\
&= \bra{\psi_1}\ket{\psi_2}
\end{align*}
Thus $U$ extends to a unitary operator $U: \bb{H} \otimes Q \lto \bb{H} \otimes Q$.

Now consider the following projective measurement on $\bb{H} \otimes Q$:
\begin{equation}
P_m := I_q \otimes \ket{m}\bra{m}
\end{equation}
Then the probability outcome $n$ occurs is:
\begin{align*}
p(n) &= \bra{\psi}\otimes\bra{\varphi}U^\dagger P_n U \ket{\psi}\otimes\ket{\varphi}\\
&= \big(U\ket{\psi}\otimes\ket{\varphi}\big)^\dagger P_n U\ket{\psi}\otimes\ket{\varphi}\\
&= \big(\sum_{m\in \call{M}}M_m\ket{\psi}\otimes\ket{m}\big)^\dagger P_n \sum_{m' \in \call{M}} M_{m'}\ket{\psi}\otimes\ket{m}\\
&= \big(\sum_{m \in \call{M}}\bra{\psi}M_m^\dagger\bra{m}\big) I_Q \otimes \ket{n}\bra{n}\sum_{m' \in \call{M}}M_{m'}\ket{\psi}\otimes\ket{m}\\
&= \big(\sum_{m \in \call{M}}\bra{\psi}M_m^\dagger \bra{m}\big)\sum_{m'\in\call{M}}M_{m'}\ket{\psi}\otimes\ket{n}\bra{n}\ket{m}\\
&=\sum_{m \in \call{M}}\bra{\psi}M_m^\dagger\bra{m} M_m \ket{\psi}\otimes \ket{n}\\
&=\sum_{m \in \call{M}}\bra{\psi}M_m^\dagger M_n\ket{\psi}\bra{m}\ket{n}\\
&=\bra{\psi}M_n^\dagger M_n \ket{\psi}
\end{align*}
\end{proof}
Now we define the \emph{evolution} of a state space:
\begin{defn}\label{def:time_evolution}
Let $\bb{H}$ be a state space. A \textbf{single step time evolution} of $\bb{H}$ is a unitary operator $U$ on $\bb{H}$. A \textbf{single step time evolution} of a state vector $\ket{\psi}$ with respect to $U$ is the pair $(\ket{\psi},U\ket{\psi})$.

An \textbf{evolution} of $\bb{H}$ is a sequence of unitary operators $(U_1,...,U_n)$ on $\bb{H}$, an \textbf{evolution} of a state vector $\ket{\psi}$ with respect to the evolution $(U_1,...,U_n)$ is the sequence $(\ket{\psi},U_1\ket{\psi},...,U_n\hdots U_1\ket{\psi})$.
\end{defn}
Lastly, the \emph{composite} of two state spaces:
\begin{defn}\label{def:composite_system}
Let $\bb{H}_1,\bb{H}_2$ be two state spaces. The \textbf{composite state space} is $\bb{H}_1 \otimes \bb{H}_2$. We often describe a state of $\bb{H}_1 \otimes \bb{H}_2$ using the terminology ``$\bb{H}_1$ is in state $\ket{\psi_1}$ and $\bb{H}_2$ is in state $\ket{\psi_2}$", this simply describes the state $\ket{\psi_1} \otimes \ket{\psi_2} \in \bb{H}_1 \otimes \bb{H}_2$.
\end{defn}

\section{The density operator}
\begin{defn}\footnote{Purity is a property of the density operator $\rho$, not of the ensemble. $\rho$ is pure iff $\rho=\ket{\psi}\bra{\psi}$, equivalently $\operatorname{Tr}(\rho^2)=1$. An ensemble of length $>1$ may still produce a pure $\rho$. See Section~\ref{sec:errata}.}
Let $\bb{H}$ be a state space. An \textbf{ensemble of states} is a finite set $\lbrace (p_1,\ket{\psi_1}),...,(p_n,\ket{\psi_n})\rbrace$ of tuples consisting of state vectors $\ket{\psi_1},...,\ket{\psi_n}$ with associated probabilities $p_1,...,p_n$, that is, $p_i \in [0,1] \subseteq \bb{R}$. We require $\sum_{i = 1}^n p_i = 1$.

If $n = 1$ then the ensemble is a \textbf{pure state}, otherwise it is a \textbf{mixed state}.

Associated to an ensemble of states is the \textbf{density operator} (sometimes \textbf{density matrix}, sometimes \textbf{state})
\begin{equation}
\sum_{i = 1}^n p_i \ket{\psi_i}\bra{\psi_i}
\end{equation}
\end{defn}
We see now that Section \ref{Sec:Math_defs} concerned itself with pure states, where we identify a state vector $\ket{\psi}$ with the operator $\ket{\psi}\bra{\psi}$ (that is, we identify the vector $\ket{\psi}$ with the projection onto this vector). We now describe how to generalise the Definitions of Section \ref{Sec:Math_defs} to the case of mixed states.

The Definition of a measurement remains identical, notice however that to a measurement $\lbrace M_m: \bb{H} \lto \bb{H}\rbrace_{m \in \call{M}}$ the associated value $p(m)$ is now a sum (NB: this is a \emph{definition} which is motivated by the law of total probability):
\begin{align*}
p(m) &= \sum_{i = 1}^n p_i p(m \mid i)\\
&=\sum_{i = 1}^n p_i \bra{\psi_i}M^\dagger_m M_m \ket{\psi_i}\\
&= \sum_{i = 1}^n p_i \operatorname{Trace}(M^\dagger_m M_m \ket{\psi_i}\bra{\psi_i})\\
&= \operatorname{Trace}(M^\dagger_m M_m \rho)
\end{align*}
where the last equality follows from linearity of the $\operatorname{Trace}$.

The resulting density operator is then \emph{defined} to be
\begin{equation}
\rho_m := \sum_{i = 1}^n p(i\mid m)\frac{M\ket{\psi_i}\bra{\psi_i}M^\dagger}{p(m \mid i)}
\end{equation}
we then use Bayes Theorem:
\begin{equation}
p(i \mid m)/p(m \mid i) = p_i/p(m)
\end{equation}
to obtain:
\begin{equation}
\rho_m = \sum_{i = 1}^n p_i\frac{M_m\ket{\psi_i}\bra{\psi_i}M_m^\dagger}{p(m)} = \sum_{i = 1}^n \frac{M_m\rho M_m^\dagger}{\operatorname{Trace}(M^\dagger_m M \rho)}
\end{equation}
We arrive at:
\begin{defn}\label{def:density_measurement}\footnote{The post-measurement density operator has the wrong operator order: it should be $\rho_m=M_m\rho M_m^\dagger/\operatorname{Tr}(M_m^\dagger M_m\rho)$, not $M_m^\dagger\rho M_m/\operatorname{Tr}(M_m^\dagger M_m\rho)$. With the correct ordering it agrees with the pure-state rule. The completeness condition is also missing $m$-subscripts. See Section~\ref{sec:errata}.}
A \textbf{measurement} on $\bb{H}$ is a family of linear operators $\lbrace M_m: \bb{H} \lto \bb{H}\rbrace_{m \in \call{M}}$ satisfying
\begin{equation}
\sum_{m \in \call{M}}M^\dagger M = I
\end{equation}
Associated to every measurement and density operator $\rho$ there is a value
\begin{equation}
p(m) = \operatorname{Trace}(M^\dagger_mM_m \rho)
\end{equation}
which is understood as the probability $p(m)$ of outcome $m$ on measurement $\lbrace M_m\rbrace_{m \in \call{M}}$.

Also, there is a \textbf{resulting density operator},:
\begin{equation}
\rho_m := \frac{M_m^\dagger \rho M_m}{\operatorname{Trace}(M^\dagger_mM_m \rho)}
\end{equation}
\end{defn}
\begin{lemma}
For a pure state density operator $\ket{\psi}\bra{\psi}$ Definitions \ref{def:density_measurement} and \ref{def:measurement} agree once $\ket{\psi}\bra{\psi}$ has been identified with $\ket{\psi}$.
\end{lemma}
Now, the Definition of time evolution:
\begin{defn}
Let $\bb{H}$ be a state space. A \textbf{single step time evolution} on $\bb{H}$ is a unitary operator $U: \bb{H} \lto \bb{H}$. A \textbf{single step time evolution} of a density operator $\rho$ with respect to $U$ is the pair $(\rho, U\rho U^\dagger)$.

An \textbf{evolution} of $\bb{H}$ is a sequence of unitary operators $(U_1,...,U_n)$ on $\bb{H}$, an \textbf{evolution} of a density operator $\rho$ with respect to the evolution $(U_1,...,U_n)$ is the sequence $(\rho, U\rho U^\dagger, ..., U_n\hdots U_1 \rho U_1^\dagger \hdots U_n^\dagger)$.\footnote{The first evolved state should be $U_1\rho U_1^\dagger$, not $U\rho U^\dagger$. See Section~\ref{sec:errata}.}
\end{defn}
Finally:
\begin{defn}
Let $\bb{H}_1,\bb{H}_2$ be two state spaces. The \textbf{composite state space }is $\bb{H}_1 \otimes \bb{H}_2$. We often describe a state of $\bb{H}_1 \otimes \bb{H}_2$ using the terminology ``$\bb{H}_1$ is in state $\rho_1$ and $\bb{H}_2$ is in state $\rho_2$",  this simply describes the state $\rho_1 \otimes \rho_2 \in \bb{H}_1 \otimes \bb{H}_2$.
\end{defn}
Let $\rho$ be an arbitrary, positive operator on a finite dimensional state space $\bb{H}$. Then since $\rho$ is positive, it is Hermitian, and thus normal. We thus obtain a spectral decomposition:
\begin{equation}
\rho = \sum_{i = 1}^n \lambda_i \ket{\psi_i}\bra{\psi_i}
\end{equation}
Thus, if $\operatorname{Trace}\rho = 1$, we have a well defined density operator with respect to the ensemble $\lbrace \lambda_i, \ket{\psi_i}\rbrace$. That every density operator is positive and has trace equal to one is an easy exercise, we arrive at:
\begin{defn}[Intrinsic definition of density operator]
Let $\bb{H}$ be finite dimensional. A \textbf{density operator} is a positive operator $\rho: \bb{H} \lto \bb{H}$ with trace equal to 1.
\end{defn}

Note that we have not described a canonical form of an ensemble, and so it is reasonable to expect different ensembles to give rise to the same density operator, indeed this is so (see \cite[Page 103]{quantum_compuiting} for an explicit example).  This leads to a natural question: when do two different ensembles give rise to the same density operator? This is answered by Proposition \ref{prop:equivalent_density_operators}, but first we make an observation, which may be skipped, to make the normalisation in the statement of Proposition \ref{prop:equivalent_density_operators} more transparent:
\begin{observation}
Let $\bb{H}$ be of dimension $n$. Let $\sum_{i = 1}^m p_i \ket{\psi}\bra{\psi}$ be a density operator associated to an ensemble. Since this is diagonalisable, we write
\begin{equation}\label{eq:canonical_attempt}
\sum_{i = 1}^m p_i \ket{\psi_i}\bra{\psi_i} = \sum_{j = 1}^n \lambda_j \ket{\varphi_j}
\end{equation}
where $\ket{\varphi_1},...,\ket{\varphi_n}$ is an orthonormal basis for $\bb{H}$. Write for each $i$:
\begin{equation}\label{eq:basis_expression}
\ket{\psi_i} = \sum_{j = 1}^n \alpha_{ij}\ket{\varphi_j}
\end{equation}
Substituting \eqref{eq:basis_expression} into \eqref{eq:canonical_attempt} to obtain:
\begin{align*}
\sum_{j = 1}^n \lambda_j \ket{\varphi_j} &= \sum_{i = 1}^m p_i \sum_{j,k = 1}^m \alpha_{ij}\overline{\alpha_{ki}} \ket{\varphi_j}\bra{\varphi_k}\\
&= \sum_{j,k = 1}^n\Big(\sum_{i = 1}^m p_i \alpha_{ij}\overline{\alpha_{ki}}\Big)\ket{\varphi_j}\bra{\varphi_k}
\end{align*}
which implies
\begin{equation}
\lambda_j = \delta_{jk}\sum_{i = 1}^m p_i \alpha_{ij}\alpha_{ij}^\dagger
\end{equation}
We mould this calculation to give the result of Proposition \ref{prop:equivalent_density_operators} by normalisation, indeed we replace \eqref{eq:canonical_attempt} with
\begin{equation}
\sqrt{p_i}\ket{\psi_i} = \sum_{j = 1}^n \alpha_{ij}\sqrt{\lambda}\ket{\varphi_i}
\end{equation}
and show that $\alpha_{ij}$ is unitary.
\end{observation}
\begin{proposition}\label{prop:equivalent_density_operators}\footnote{This statement is mathematically incoherent: the notation $\{p_i\ket{i}\bra{i}\}$ is not standard ensemble notation; indices, kets/bras, and quantifiers are conflated; and the proof omits bras in spectral decompositions. See Section~\ref{sec:errata} for the standard unitary-freedom theorem for ensembles.}
Let $\lbrace p_i \ket{i}\bra{i}\rbrace ,\lbrace  q_j \ket{j}\bra{j}\rbrace$ be two ensembles and denote by $\rho_i,\rho_j$ their respective induced density operators. We Then $\rho_i = \rho_j$ if and only if for each $i = 1,...,m$ there exists a unitary matrix $u^i_{jk}$ such that $$\sqrt{p_i}\ket{\psi_i} = \sum_{j = 1}^n \alpha_{ij}\sqrt{q_j}\ket{\varphi_j}$$ with $\alpha_{ij}$ a unitary matrix.
\end{proposition}
\begin{proof}
Write $\sum_{j = 1}^n \lambda_j \ket{\varphi_j}$ for the diagonalisation of $\rho_i$ and write
\begin{equation}
\sqrt{p_i}\ket{\psi_i} = \sum_{j = 1}^n \sqrt{\lambda_j}\ket{\varphi_j}
\end{equation}
Then
\begin{align*}
\sum_{i = 1}^m p_i \ket{\psi_i}\bra{\psi_i} &= \sum_{i = 1}^m \sqrt{p_i}\ket{\psi_i}\sqrt{p_i}\bra{\psi_i}\\
&= \sum_{i = 1}^m  \sum_{j,k = 1}^n\alpha_{ij}\overline{\alpha_{kj}}\sqrt{\lambda_j}\ket{\varphi_j}\sqrt{\lambda_k}\bra{\varphi_{k}}\\
&= \sum_{j,k = 1}^n \Big( \alpha_{ij}\overline{\alpha_{kj}}\Big)\sqrt{\lambda_j}\ket{\varphi_j}\sqrt{\lambda_k}\bra{\varphi_{k}}
\end{align*}
We have by assumption definition that $\sum_{i = 1}^m p_i \ket{\psi_i}\bra{\psi_i}  = \sum_{i = 1}^n \lambda_i \ket{\varphi_i}\bra{\varphi_i} = \sum_{i = 1}^n \sqrt{\lambda_i}\ket{\varphi_i}\sqrt{\lambda_i}\bra{\varphi_i}$. It follows that $\alpha_{ij}$ is unitary.  Applying the same logic to $\rho_j$ we see that both $\rho_i,\rho_j$ are related to $\sum_{j = 1}^n \lambda_j \ket{\varphi_j}\bra{\varphi_i}$ and so we are done.
\end{proof}



\subsection{Partial trace}

For an introduction to the partial trace operator, see \cite{CommutativeAlgebra}
\begin{example}
We calculate the partial trace of the operator
\begin{equation}
\rho := \Big(\frac{\ket{00} + \ket{11}}{\sqrt{2}}\Big)\Big(\frac{\bra{00} + \bra{11}}{\sqrt{2}}\Big) = \frac{\ket{00}\bra{00} + \ket{00}\bra{11} + \ket{11}\bra{00} + \ket{11}\bra{11}}{2}
\end{equation}
We denote the first copy of $\bb{C}^2$ simply by $1$ and the second copy of $\bb{C}^2$ by $2$.

First consider $\operatorname{Trace}_2(\ket{00}\bra{00})$. We have, where we write $E_{ij}: \bb{C}^2 \lto \bb{C}^2$ for the linear map which maps the $i^{\text{th}}$ basis vector to the $j^{\text{th}}$ basis vector, and similarly for $F_{ij}$ (just applied to the second copy of $\bb{C}^2$)
\begin{equation}
\ket{00}\bra{00} = E_{00} \otimes F_{00}
\end{equation}
and so
\begin{equation}
\operatorname{Trace}_2(\ket{00}\bra{00}) = \operatorname{Trace}(F_{00})E_{00} = \ket{0}\bra{0}
\end{equation}
Similarly,
\begin{equation}
\ket{11}\bra{00} = E_{01}\otimes F_{01},\quad \ket{00}\bra{11} = E_{10} \otimes F_{10},\quad \ket{11}\bra{11} = E_{11} \otimes F_{11}
\end{equation}
and so
\begin{align}
\operatorname{Trace}_2(\ket{11}\bra{00}) = \operatorname{Trace}(F_{00})E_{11} &= 0\\
\operatorname{Trace}_2(\ket{11}\bra{00}) = \operatorname{Trace}(F_{00})E_{11} &= 0\\
\operatorname{Trace}_2(\ket{11}\bra{11}) = \operatorname{Trace}(F_{11})E_{11} &= \ket{1}\bra{1}
\end{align}\footnote{The $\ket{11}\bra{00}$ line is duplicated and the trace entries are mislabelled. The correct rows are $\operatorname{Tr}_2(\ket{11}\bra{00})=\operatorname{Tr}(F_{01})E_{01}=0$, $\operatorname{Tr}_2(\ket{00}\bra{11})=\operatorname{Tr}(F_{10})E_{10}=0$, $\operatorname{Tr}_2(\ket{11}\bra{11})=\operatorname{Tr}(F_{11})E_{11}=\ket{1}\bra{1}$. See Section~\ref{sec:errata}.}
\end{example}
we thus have
\begin{equation}
\operatorname{Trace}_2(\rho) = \frac{\ket{0}\bra{0} + \ket{1}\bra{1}}{2} = I/2
\end{equation}
We present another way defining the partial trace:
\begin{observation}\footnote{The standard tensor-Hom adjunction is $\operatorname{Hom}(V\otimes U,W)\cong\operatorname{Hom}(U,V^*\otimes W)$. The right adjoint of $V\otimes-$ is $V^*\otimes-$. See Section~\ref{sec:errata}.}
Let $V$ be a finite dimensional vector space. The functor $V \otimes {\und{0.2}}$ admits a right adjoint which is given by the functor $\und{0.2} \otimes V^\ast$:
\begin{equation}
\operatorname{Hom}(V \otimes U, W) \cong \operatorname{Hom}(V, W \otimes V^\ast)
\end{equation}
\end{observation}

\subsection{Quantum operations}
\begin{defn}
An \textbf{open quantum system} is the tensor product of two state spaces $\bb{H}_p \otimes \bb{H}_e$ where $\bb{H}_p$ is the \textbf{principal system} and $\bb{H}_e$ is the \textbf{environment}.
\end{defn}
\begin{defn}
The \textbf{time evolution of an open quantum system} is that of a Definition \ref{def:time_evolution} where an open quantum system is thought of as a composite system (Definition \ref{def:composite_system}).

Given an open quantum system $\bb{H}_p \otimes \bb{H}_e$, and a state $\rho \otimes \rho_e$ a \textbf{quantum operation} is an operator $\call{E}: \bb{H}_p \lto \bb{H}_p$\footnote{Type-incorrect: a quantum operation is a linear map on \emph{operators}, $\mathcal{E}:L(\mathbb{H}_p)\to L(\mathbb{H}_p)$, not on the Hilbert space itself. Trace-preserving CP such maps are called quantum channels. See Section~\ref{sec:errata}.} given by the following formula, where $U$ is a unitary operator $U$ on the entire open quantum system $\bb{H}_p \otimes \bb{H}_e$:
\begin{equation}\label{eq:quant_op_general_form}
\call{E}(\rho) = \operatorname{Trace}_{\bb{H}_e}(U (\rho \otimes \rho_e) U^\dagger)
\end{equation}
\end{defn}
\begin{remark}\label{rmk:operator_sum_derivation}
Let $\bb{H}_p \otimes \bb{H}_e$ be an open quantum system and let $\ket{1},...,\ket{n}$ be a basis for $\bb{H}_e$. We think of $\ket{1},...,\ket{n}$ as operators $\bb{C} \lto V$ and write $\operatorname{id} \otimes \ket{i}$ for the composite $\bb{H}_p \lto \bb{H}_p \otimes \bb{C} \lto \bb{H}_p \otimes \bb{H}_e$. Also, we denote the multiplication map $\bb{H}_p \otimes k \lto \bb{H}_p$ defined by $w \otimes x \longmapsto xw$ by $\operatorname{Mult}$. We have for any operator $f: \bb{H}_p \otimes \bb{H}_e \lto \bb{H}_p \otimes \bb{H}_e$ that
\begin{equation}\label{eq:partial_trace_exp_form}
\operatorname{Trace}_{\bb{H}_e}f = \operatorname{Mult}\Big(\sum_{i = 1}^n (\operatorname{id} \otimes \bra{i})f(\operatorname{id} \otimes \ket{i})\Big)
\end{equation}
See \cite{CommutativeAlgebra} for more details.
\end{remark}
We let $\ket{1},...,\ket{n}$ be a basis for $\bb{H}_e$ and use \eqref{eq:partial_trace_exp_form} to rewrite \eqref{eq:quant_op_general_form} in the special case where $\rho_e = \ket{e_0}\bra{e_0}$.\footnote{The derivation conflates the summation index $i$ with the fixed environment state, writing $\rho\otimes\ket{i}\bra{i}$ where $\rho\otimes\ket{e_0}\bra{e_0}$ is intended. See Section~\ref{sec:errata}.}We have
\begin{align}
\call{E}(\rho) &= \operatorname{Trace}_{\bb{H}_e}(U (\rho \otimes \ket{i}\bra{i})U^\dagger)\\
&= \operatorname{Mult}\Big(\sum_{i = 1}^n (\operatorname{id} \otimes \bra{i}) (U (\rho \otimes \ket{i}\bra{i})U^\dagger)(\operatorname{id} \otimes \ket{i})\Big)\label{eq:trace_mult}
\end{align}
Now we make the observation that
\begin{align}
\rho \otimes \ket{i}\bra{i} &= (\rho \otimes \operatorname{id})(\operatorname{id} \otimes \ket{e_0})(\operatorname{id} \otimes \bra{e_0})\\
&= (\operatorname{id} \otimes \ket{e_0})(\rho \otimes \operatorname{id})(\operatorname{id} \otimes \bra{e_0})\label{eq:tensor_trick}
\end{align}
Substituting \eqref{eq:tensor_trick} into \eqref{eq:trace_mult} we obtain:
\begin{align*}
&\operatorname{Mult}\Big(\sum_{i = 1}^n (\operatorname{id} \otimes \bra{i}) (U (\rho \otimes \ket{i}\bra{i})U^\dagger)(\operatorname{id} \otimes \ket{i})\Big)\\
= &\operatorname{Mult}\Big(\sum_{i = 1}^n (\operatorname{id} \otimes \bra{i}) U (\operatorname{id} \otimes \ket{e_0})(\rho \otimes \operatorname{id})(\operatorname{id} \otimes \bra{e_0})U^\dagger(\operatorname{id} \otimes \ket{i})\Big)\\
= &\sum_{i = 1}^n \operatorname{Mult}\Big((\operatorname{id} \otimes \bra{i}) U (\operatorname{id} \otimes \ket{e_0})\Big)\rho\operatorname{Mult}\Big((\operatorname{id} \otimes \bra{e_0})U^\dagger(\operatorname{id} \otimes \ket{i})\Big)\\
= &\sum_{i = 1}^n E_i \rho E_i^\dagger
\end{align*}
where $E_i = \operatorname{Mult}\Big((\operatorname{id} \otimes \bra{i}) (U ((\operatorname{id} \otimes \ket{e_0})\Big)$.

We thus have a second Definition of a quantum operation:
\begin{lemma}\label{lem:operator_sum}\footnote{A quantum operation is not itself a finite set of operators; it is a linear map admitting an operator-sum representation. The Kraus operators are not unique. See Section~\ref{sec:errata}.}
Given an open quantum system $\bb{H}_p \otimes \bb{H}_e$, a quantum operation $\call{E}$ is a finite set of operators $\lbrace E_1,...,E_n\rbrace$ where each $E_i$ is an operator on the principal system $\bb{H}_p$ and the collection $\call{E}$ satisfies the \textbf{completeness condition}, that
\begin{equation}
\sum_{i = 1}^n E_i^\dagger E_i = I
\end{equation}
\end{lemma}
Notice that Definition \ref{lem:operator_sum} does not include the data of $U$. Indeed, given just the principle system and the operator elements $\lbrace E_1,...,E_n\rbrace$, an environment $\bb{H}_e$ and a unitary operator $U: \bb{H}_p \otimes \bb{H}_e \lto \bb{H}_p \otimes \bb{H}_e$ can be defined which give rise to $\lbrace E_1,...,E_n\rbrace$ following the calculation of Remark \ref{rmk:operator_sum_derivation}. For details see \cite[\S 8.2.3 page 365]{quantum_computing}.

\section{Quantum error correction codes}
Throughout this Section, we denote by $\bb{H}$ a single qubit $\bb{C}^2$.
\begin{defn}\footnote{Basis convention is reversed: the standard convention (consistent with $Z=\operatorname{diag}(1,-1)$, used throughout the rest of the document) is $\ket{0}=(1,0)^T$ and $\ket{1}=(0,1)^T$. See Section~\ref{sec:errata}.}
Let $\ket{0}, \ket{1}$ respectively denote the standard basis vectors $(0,1),(1,0)$ of the $\bb{C}$-vector space $\bb{C}^2$. If $a\ket{0} + b\ket{1}$ is a state and both $a$ and $b$ are non-zero then the state is a \textbf{superposition}.
\end{defn}

\begin{defn}
A \textbf{quantum gate} is a unitary matrix in $M_2(\bb{C})$, and is thought of as an operator $\bb{H} \lto \bb{H}$ (written with respect to the standard basis $\ket{0},\ket{1}$).
\end{defn}
The game of quantum computation is to build a computer out of quantum gates. For some basic examples see \cite[\S 4]{quantum_computing}.

Of all the quantum gates, there are a few which are ubiquitous in the theory:
\begin{defn}
We denote the following operators:
\begin{equation}
X :=
\begin{pmatrix}
0 &1\\
1 &0
\end{pmatrix}
\quad
Y :=
\begin{pmatrix}
0 & -i\\
i & 0
\end{pmatrix}
\quad
Z :=
\begin{pmatrix}
1 & 0\\
0 & -1
\end{pmatrix}
\quad
H :=
\frac{1}{\sqrt{2}}
\begin{pmatrix}
1 & 1\\
1 & -1
\end{pmatrix}
\end{equation}
The matrices $X,Y,Z$ are the \textbf{Pauli matrices}, and $H$ is the \textbf{Hadamard matrix}.
\end{defn}
We make the passing observation that all of $X,Y,Z,H$ square to the identity matrix. The basis vectors
\begin{equation}
H\ket{0} = \frac{1}{\sqrt{2}}(\ket{0} + \ket{1}),\quad H\ket{1} = \frac{1}{\sqrt{2}}(\ket{0} - \ket{1})
\end{equation}
are the \textbf{Bell states}\footnote{These are not Bell states. They are the \emph{Hadamard-basis} (or $X$-basis) states $\ket{\pm}$. Bell states are two-qubit entangled states such as $(\ket{00}+\ket{11})/\sqrt{2}$. See Section~\ref{sec:errata}.} and are denoted $\ket{+},\ket{-}$ respectively. Notice that as already stated, $H^2 = I$, so $H\ket{+} = \ket{0}$ and $H\ket{-} = \ket{1}$.
\begin{defn}
The standard basis $\ket{0},\ket{1}$ of $\bb{H}$ induces a basis of $\bb{H}^{\otimes n}$, we denote $\ket{0} \otimes \hdots \otimes \ket{0}$ by $\ket{0\hdots 0}$, etc.
\end{defn}

\begin{notation}
Given a collection of Pauli matrices $W_{i_1},...,W_{i_m} \in \lbrace X,Y,Z\rbrace$ where $0 < i_1 < \hdots < i_m < n$ we denote by $W_{i_1}...W_{i_m}$ the operator on $\bb{H}^{\otimes n}$ given by the tensor product with $W_{i_j}$ in the $i_j^{\text{th}}$ slot and with the identity operator in all other slots. For example, the operator $Z_1$ on $\bb{H}^{\otimes 3}$ is the operator $Z \otimes I \otimes I$.
\end{notation}


\begin{defn}
A \textbf{message} is a state $\ket{\psi} \in \bb{H}^{\otimes n}$, for some $n$. An \textbf{error} is a pair of states $(\ket{\varphi},\ket{\psi})$ where $\ket{\varphi},\ket{\psi} \in \bb{H}^{\otimes n}$ for some $n$, note that an error may be such that $\ket{\varphi} = \ket{\psi})$.  The message $\ket{\varphi}$ is the \textbf{intended message} and $\ket{\psi}$ is the \textbf{received message}. 
\end{defn}
\begin{defn}\footnote{A quantum encoding should be a linear isometry, not merely an injective map; otherwise superpositions and inner products need not be preserved. See Section~\ref{sec:errata}.}
An \textbf{$n$-encoding of a single state} (sometimes just an \textbf{encoding}) is an injective map $\iota: \bb{H} \lto \bb{H}^{\otimes n}$. An \textbf{$n$-encoding of a message} $\ket{m} \in \bb{H}^{\otimes k}$ is an $n$-encoding $\iota$ along with a message $\ket{m} \in \bb{H}^{\otimes nk}$ for which there exists $\ket{m'} \in \bb{H}^{\otimes k}$ satisfying $\iota^{\otimes k}\ket{m'} = \ket{m}$.
\end{defn}
\begin{example}\label{ex:bit_flip_coding}
Consider the encoding
\begin{align}\label{enc:bit_flip}
\operatorname{BitFlip}: \bb{H} &\lto \bb{H}^{\otimes 3}\\
\ket{0} &\longmapsto \ket{000}\\
\ket{1} &\longmapsto \ket{111}
\end{align}
then an encoding of a message with respect to this encoding might be $\ket{000111000}$, but could not be $\ket{000111001}$. We call Encoding \ref{enc:bit_flip} the \textbf{bit flip encoding}, for reasons made clear by Algorithm \ref{alg:bit_flip_correction} below. As another example, we consider the \textbf{phase flip encoding}:
\begin{align*}
\operatorname{PhaseFlip}: \bb{H} &\lto \bb{H}^{\otimes 3}\\
\ket{0} &\longmapsto \ket{+++}\\
\ket{1} &\longmapsto \ket{---}
\end{align*}
\end{example}
\begin{defn}
Let $\bb{H}$ be a state space. A \textbf{quantum error correction code (QECC)} for $\bb{H}$ is a subspace $C \subseteq \bb{H}$.
\end{defn}
\begin{example}
Continuing with Example \ref{ex:bit_flip_coding}, two QECCs are given by
\begin{equation}
\operatorname{im}(\operatorname{BitFlip}),\qquad\text{and}\qquad\operatorname{im}(\operatorname{PhaseFlip})
\end{equation}
\end{example}
\begin{defn}
Given a state space $\bb{H}$ and a QECC $\bb{W} \subseteq \bb{H}$, a \textbf{correctable set of errors} is a quantum operation $\scr{E}$ for which there exists a trace preserving quantum operation $\call{R}$ satisfying:
\begin{equation}
\call{R}(\scr{E}(\rho))\propto \rho
\end{equation}
\end{defn}
\begin{defn}
A \textbf{bit flip error} is an error $(\ket{\varphi},\ket{\psi})$ where $\ket{\varphi}$ is an encoding of a message with respect to the encoding $\operatorname{BitFlip}^{\otimes m}$ for some $m$, such that $X_i\ket{\varphi} =\ket{\psi}$ for some $i$.

A \textbf{phase flip error} is an error  $(\ket{\varphi},\ket{\psi})$ where $\ket{\varphi}$ is an encoding of a message with respect to the the encoding $\operatorname{PhaseFlip}^{\otimes m}$, such that $Z_i\ket{\varphi} =\ket{\psi}$ for some $i$.
\end{defn}
Let $(\ket{\varphi},\ket{\psi})$ be a bit flip error.  The following algorithm takes as input $\ket{\psi}$ and reconstructs $\ket{\varphi}$:
\begin{algorithm}[Bit flip correction]\label{alg:bit_flip_correction}\footnote{The expression $\bra{\psi}Z_1Z_2\ket{\psi}$ is the expectation value, not the projective measurement. The projective measurement of $Z_1Z_2$ uses projectors $P_\pm=(I\pm Z_1Z_2)/2$. Code states are eigenvectors of $Z_1Z_2,Z_2Z_3$, so the syndrome is revealed without disturbing the state. Applying $Z_1Z_2$ as a unitary (as the proof does) is a different operation. See Section~\ref{sec:errata}.}
Input: a received message $\ket{\psi}$,
\begin{enumerate}
\item perform the following projective measurements:
\begin{equation}
\bra{\psi} Z_1 Z_2 \ket{\psi}\text{ with resulting state }\ket{\psi'},
\end{equation}
followed by
\begin{equation}
\bra{\psi'} Z_2 Z_3 \ket{\psi'}
\end{equation}
let $(r_1,r_2)$ be the pair of results from these measurements.
\item It will be shown that $r_1,r_2 \in \lbrace 1,-1\rbrace$, and the resulting state of the second measurement is $\ket{\psi}$.
\item Now retrieve $\ket{\varphi}$ based on the values of $r_1,r_2$:
\begin{itemize}
\item if $(r_1, r_2) = (1,1)$, return $\ket{\psi}$,
\item if $(r_1,r_2) = (-1,1)$, return $X_1 \ket{\psi}$,
\item if $(r_1,r_2) = (1,-1)$, return $X_3 \ket{\psi}$,
\item if $(r_1,r_2) = (-1,-1)$, return $X_2 \ket{\psi}$
\end{itemize}
\end{enumerate}
\end{algorithm}
We now prove correctness of Algorithm \ref{alg:bit_flip_correction}:
\begin{proof}
It will be helpful to first notice:
\begin{align*}
Z_1Z_2\ket{000} &= \ket{000} & Z_1Z_2\ket{001} &= \ket{001}\\
Z_1Z_2\ket{010} &= -\ket{010} & Z_1Z_2\ket{011} &= -\ket{011}\\
Z_1Z_2\ket{100} &= -\ket{100} & Z_1Z_2\ket{101} &= -\ket{101}\\
Z_1Z_2\ket{110} &= \ket{110} &Z_1Z_2 \ket{111} &= \ket{111}
\end{align*}
Let $\ket{\psi}:= a\ket{010} + b\ket{101}$ be a state, ie, an element of $\bb{H}^{\otimes 3}$. We perform the measurement $Z_1Z_2$ followed by $Z_2Z_3$:
\begin{align*}
\bra{\psi}Z_1Z_2\ket{\psi} &= (a\bra{010} + b\bra{101})Z_1Z_2(a\ket{010} + b\ket{101})\\
&= (a\bra{010} + b\bra{101})(-a\ket{010} - b\ket{101})\\
&= -a^2 - b^2 = -1\footnote{For $a,b\in\mathbb{C}$ this should read $-|a|^2-|b|^2=-1$, with bras $\overline{a}\bra{010}+\overline{b}\bra{101}$. See Section~\ref{sec:errata}.}
\end{align*}
and
\begin{align*}
\bra{\psi}Z_2Z_3\ket{\psi} &= (a\bra{010} + b\bra{101})Z_1Z_2(a\ket{010} + b\ket{101})\\
&= (a\bra{010} + b\bra{101})(-a\ket{010} - b\ket{101})\\
&= -a^2 - b^2 = -1
\end{align*}
We can infer from the fact that both of these came out as $-1$ that it was the second bit which was flipped, and so we can correct this. However, what is the impact of this measurement on the state? Again we calculate:
\begin{align*}
Z_1Z_2(a\ket{010} + b\ket{101}) &= Z_1(-a\ket{010} + b\ket{101})\\
&= -a\ket{010} - b\ket{101}
\end{align*}
and
\begin{align*}
Z_2Z_3(-a\ket{010} - b\ket{101}) &= Z_2(-a\ket{010} + b\ket{101})\\
&= a\ket{010} + b\ket{101}
\end{align*}
and so the measurements (in the end) did not impact our state.
\end{proof}
Later, using the theory of \emph{stabiliser codes}, we will show that in fact single bit flip errors form the full set of correctable errors using $Z_1Z_2,Z_1Z_3,Z_2Z_3$. We can rephrase what we have demonstrated in the language introduced at the beginning of this Section:
\begin{example}\footnote{The correctable errors are $\{I,X_1,X_2,X_3\}$, not $\{Z_1Z_2,Z_1Z_3,Z_2Z_3\}$. The latter are syndrome observables. See Section~\ref{sec:errata}.}
Given the state space $\bb{H} = (\bb{C}^2)^{\otimes 3}$ and QECC $\bb{W} \subseteq \bb{H}$ given by $\bb{W} = \operatorname{im}\operatorname{BitFlip}$, a correctable set of errors is $\scr{E} = \lbrace Z_1Z_2,Z_1Z_3,Z_2Z_3\rbrace$. The corresponding trace preserving quantum operation $\call{R}$ is given by
\begin{equation}
\call{R}(\rho) = 
\end{equation}
\end{example}

Let $(\ket{\varphi},\ket{\psi})$ be a phase flip error. The following algorithm takes as input $\ket{\psi}$ and reconstructs $\ket{\varphi}$:
\begin{algorithm}[Phase flip correction]\footnote{Same measurement-versus-application confusion as the bit-flip case: projective measurements use $(I\pm X_1X_2)/2$ and $(I\pm X_2X_3)/2$. See Section~\ref{sec:errata}.}
Input: a received message $\ket{\psi}$:
\begin{enumerate}
\item perform the following projective measurements:
\begin{equation}
\bra{\psi} X_1 X_2 \ket{\psi} \text{ with resulting state }\ket{\psi'}
\end{equation}
followed by
\begin{equation}
\bra{\psi'} X_2 X_3 \ket{\psi'}
\end{equation}
let $(r_1,r_2)$ be the pair of results from these measurements.
\item It will be shown that $r_1, r_2 \in \lbrace 1, -1\rbrace$ and the resulting state of the second measurement is $\ket{\psi}$.
\item Now retrieve $\ket{\varphi}$ based on the values of $r_1,r_2$:
\begin{enumerate}
\item if $(r_1,r_2) = (1,1)$, return $\ket{\psi}$,
\item if $(r_1,r_2) = (-1,1)$, return $Z_1 \ket{\psi}$
\item if $(r_1,r_2) = (1,-1)$, return $Z_3 \ket{\psi}$,
\item if $(r_1,r_2) = (-1,-1)$, return $Z_2 \ket{\psi}$
\end{enumerate}
\end{enumerate}
\end{algorithm}
\begin{proof}
In fact all our work is already done. We simply note that $Z\ket{+} = \ket{-}, Z\ket{-} = \ket{+}$ (and so phase flip acts like bit flip for $\ket{+},\ket{-})$, and that  in general
\begin{equation}
H^{\otimes n}Z_{i_1}\hdots Z_{i_j}H^{\otimes n} = X_{i_1}\hdots X_{i_j}
\end{equation}
The result then follows from the proof of correctness for Algorithm \ref{alg:bit_flip_correction}.
\end{proof}
What if we wanted to correct an error where we knew the received message corresponded to the intended message by either a bit flip \emph{or} a phase flip. This can be done by combining the two approaches above. Define the following encoding:
\begin{defn}\footnote{This composition is not type-correct: $\operatorname{PhaseFlip}:\mathbb{H}\to\mathbb{H}^{\otimes 3}$ produces three qubits, but $\operatorname{BitFlip}:\mathbb{H}\to\mathbb{H}^{\otimes 3}$ accepts a single qubit. The correct composition is $\operatorname{Shor}=\operatorname{BitFlip}^{\otimes 3}\circ\operatorname{PhaseFlip}$. See Section~\ref{sec:errata}.}
The \textbf{Shor encoding} is:
\begin{equation}
\operatorname{Shor}: \bb{H} \lto \bb{H}^{\otimes 9}
\end{equation}
where
\begin{equation}
\operatorname{Shor}(\ket{\psi}) = \operatorname{BitFlip} \circ \operatorname{PhaseFlip}\ket{\psi}
\end{equation}
\end{defn}
\begin{algorithm}\footnote{The listed $Z$-type measurements are not the Shor stabilisers. Cross-block checks $Z_3Z_4$ and $Z_6Z_7$ disturb the encoded information. The correct $Z$-type generators are $Z_1Z_2,Z_2Z_3,Z_4Z_5,Z_5Z_6,Z_7Z_8,Z_8Z_9$ (within-block only). The recovery table below has length-9 rows where length-8 was claimed, repeats rows, and mixes recovery rules. See Section~\ref{sec:errata} for the standard formulation.}
On input $\ket{\psi}$:
\begin{enumerate}
\item Perform the following projective measurements:
\begin{align*}
&\bra{\psi}Z_1Z_2\ket{\psi}\text{ with resulting state }\ket{\psi'}\\
&\bra{\psi'}Z_2Z_3\ket{\psi'}\text{ with resulting state }\ket{\psi}\\
&\bra{\psi}Z_3Z_4\ket{\psi}\text{ with resulting state }\ket{\psi''}\\
&\bra{\psi'}Z_4Z_5\ket{\psi''}\text{ with resulting state }\ket{\psi}\\
&\bra{\psi'}Z_5Z_6\ket{\psi}\text{ with resulting state }\ket{\psi'''}\\
&\bra{\psi'}Z_6Z_7\ket{\psi'''}\text{ with resulting state }\ket{\psi}\\
&\bra{\psi'}Z_7Z_8\ket{\psi}\text{ with resulting state }\ket{\psi''''}\\
&\bra{\psi'}Z_8Z_9\ket{\psi''''}\text{ with resulting state }\ket{\psi}
\end{align*}
let $(r_1,r_2,r_3,r_4,r_5,r_6,r_7,r_8) \in \bb{Z}_2^8$ be the results from these measurements. Notice that there are only three possibilities, all entries are $1$, exactly one entry is $-1$ in which case it is either $r_1$ or $r_8$ (with the rest equal to $1$) or exactly two values are $-1$ and the rest are $1$ in which case the two $-1$ entries are neighbours.
\item Then perform the following measurements:
\begin{align*}
&\bra{\psi}X_1X_2X_3X_4X_5X_6\ket{\psi}\text{ with resulting state }\ket{\psi'}\\
&\bra{\psi'}X_4X_5X_6X_7X_8X_9\ket{\psi'}\text{ with resulting state }\ket{\psi}
\end{align*}
let $(s_1,s_2) \in \bb{Z}_2^2$ be the result of these measurements.
\end{enumerate}
\item Now retrieve $\ket{\varphi}$ based on the values:
\begin{center}
\begin{tabular}{|c|c|c|}
\hline
$(r_1,r_2,r_3,r_4,r_5,r_6,r_7,r_8)$ & $(s_1,s_2)$ & Return\\
\hline
$(1,1,1,1,1,1,1,1)$ & $(1,1)$ & $\ket{\psi}$\\
\hline
$(-1,1,1,1,1,1,1,1,1)$ & $(1,1)$ & $X_1\ket{\psi}$\\
\hline
$(-1,-1,1,1,1,1,1,1,1)$ & $(1,1)$ & $X_2\ket{\psi}$\\
\hline
$\vdots$ & $\vdots$ & $\vdots$\\
\hline
$(1,1,1,1,1,1,1,-1,-1)$ & $(1,1)$ & $X_8\ket{\psi}$\\
\hline
$(1,1,1,1,1,1,1,1,-1)$ & $(1,1)$ & $X_9\ket{\psi}$\\
\hline
$(1,1,1,1,1,1,1,1)$ & $(-1,1)$ & $\ket{\psi}$\\
\hline
$(-1,1,1,1,1,1,1,1,1)$ & $(-1,1)$ & $Z_1Z_2Z_3X_1\ket{\psi}$\\
\hline
$(-1,-1,1,1,1,1,1,1,1)$ & $(-1,1)$ & $Z_1Z_2Z_3X_2\ket{\psi}$\\
\hline
$\vdots$ & $\vdots$ & $\vdots$\\
\hline
$(1,1,1,1,1,1,1,-1,-1)$ & $(-1,-1)$ & $Z_4Z_5Z_6X_8\ket{\psi}$\\
\hline
$(1,1,1,1,1,1,1,1,-1)$ & $(-1,-1)$ & $Z_4Z_5Z_6X_9\ket{\psi}$\\
\hline
$(-1,1,1,1,1,1,1,1,1)$ & $(-1,-1)$ & $Z_4Z_5Z_6X_1\ket{\psi}$\\
\hline
$(-1,-1,1,1,1,1,1,1,1)$ & $(-1,-1)$ & $Z_4Z_5Z_6X_2\ket{\psi}$\\
\hline
$\vdots$ & $\vdots$ & $\vdots$\\
\hline
$(1,1,1,1,1,1,1,-1,-1)$ & $(-1,-1)$ & $Z_4Z_5Z_6X_8\ket{\psi}$\\
\hline
$(1,1,1,1,1,1,1,1,-1)$ & $(-1,-1)$ & $Z_4Z_5Z_6X_9\ket{\psi}$\\
\hline
$(1,1,1,1,1,1,1,-1,-1)$ & $(-1,-1)$ & $Z_4Z_5Z_6X_8\ket{\psi}$\\
\hline
$(1,1,1,1,1,1,1,1,-1)$ & $(-1,-1)$ & $Z_4Z_5Z_6X_9\ket{\psi}$\\
\hline
$(-1,1,1,1,1,1,1,1,1)$ & $(-1,1)$ & $Z_7Z_8Z_9X_1\ket{\psi}$\\
\hline
$(-1,-1,1,1,1,1,1,1,1)$ & $(-1,1)$ & $Z_7Z_8Z_9X_2\ket{\psi}$\\
\hline
$\vdots$ & $\vdots$ & $\vdots$\\
\hline
$(1,1,1,1,1,1,1,-1,-1)$ & $(-1,1)$ & $Z_7Z_8Z_9X_8\ket{\psi}$\\
\hline
$(1,1,1,1,1,1,1,1,-1)$ & $(-1,1)$ & $Z_7Z_8Z_9X_9\ket{\psi}$\\
\hline
\end{tabular}
\end{center}
\end{algorithm}
\begin{proof}\footnote{``Obvious'' is not enough given that the listed measurements are not the standard Shor stabilisers. The combined-error step is correct only up to a global phase, which deserves explicit treatment via Pauli (anti)commutation. See Section~\ref{sec:errata}.}
That the Shor algorithm correct bit flip errors is obvious.

Now assume a single phase flip error has occurred, and no bit flip error has occurred. The core observation is the commutativity of the following diagrams for any $\ket{v} \in \lbrace \ket{0},\ket{1}\rbrace$ (where we think of any ``ket" vector $\ket{v}$ as a map $k \lto \ket{v}$) we have written $\operatorname{BF}$ for $\operatorname{BitFlip}$:
\begin{equation}
\begin{tikzcd}[column sep = huge, row sep = huge]
k\arrow[dr,swap,"{\ket{vvv}}"]\arrow[r,"{\ket{v}}"] & \bb{H}\arrow[r,"X"]\arrow[d,"{\operatorname{BF}}"] & \bb{H}\arrow[r,"{\bra{v}}"]\arrow[d,"{\operatorname{BF}}"] & k\\
& \bb{H}^{\otimes 3}\arrow[r,"{X\otimes XZ \otimes X}"] & \bb{H}^{\otimes 3}\arrow[ur,swap,"{\bra{vvv}}"]
\end{tikzcd}
\end{equation}
A similar Diagram but with $X\otimes XZ \otimes X$ replaced by $XZ\otimes X \otimes X$ or by $X\otimes X \otimes XZ$ also commutes. Thus, if $\ket{\psi} = Z_i\ket{\varphi}$, defining
\begin{equation}
s(i) =
\begin{cases}
1, & i = 1,2,3\\
2, & i=4,5,6\\
3, & i=7,8,9
\end{cases}
\end{equation}
we have
\begin{equation}
\bra{m}\operatorname{BF}^{\otimes 3\dagger}Z_i^\dagger X_1X_2X_3X_4X_5X_6 Z_i \operatorname{BF}^{\otimes 3}\ket{m} = \bra{m}Z_{s(i)}^\dagger X_1X_2Z_{s(i)}\ket{m}
\end{equation}
and so we can treat each ``block" of three states as a single state, so we know how to interpret the measurements $(s_1,s_2)$. The last observation to make is commutativity of the following Diagram for all $i=1,2,3$
\begin{equation}
\begin{tikzcd}[column sep = huge, row sep = huge]
\bb{H}^{\otimes 3}\arrow[r,"{Z_i}"]\arrow[d,"{\operatorname{BF}}"] & \bb{H}^{\otimes 3}\arrow[d,"{\operatorname{BF}}"]\\
\bb{H}^{\otimes 9}\arrow[r,"{Z_{3i-2}Z_{3i-1}Z_{3i}}"] & \bb{H}^{\otimes 9}
\end{tikzcd}
\end{equation}
Now say a combination of a bit flip and a phase flip error occurred. That is, say $\ket{\psi} = Z_iX_j\ket{\varphi}$. The error correction will first correct the bit flip which reduces to the previous case. In other words, $X_j^2 \ket{\varphi} = \ket{\varphi}$ is in the image of $\operatorname{BitFlip}$.
\end{proof}

\section{The stabiliser formalism}
\begin{defn}\footnote{Cleaner statement: $G_n=\{\alpha P_1\otimes\cdots\otimes P_n : \alpha\in\{\pm 1,\pm i\},\, P_j\in\{I,X,Y,Z\}\}$. The body's binary representation says ``$g^j=X$ iff $x_j=1$'', which fails for $g^j=Y$ (which also has $x_j=1$); the correct conditions are $x_j=1\iff g^j\in\{X,Y\}$ and $z_j=1\iff g^j\in\{Y,Z\}$. There is also an index typo $x_i$ for $x_j$. See Section~\ref{sec:errata}.}
The \textbf{Pauli Group}, denoted $G_n$, is generated by of all operators $\pm I,  i I, X_j, Y_j, Z_j$ for $j = 1,...,n$.

Denote by $\scr{X}$ denote the following set of operators
\begin{equation}
\scr{X} := \lbrace I,X,Y,Z\rbrace
\end{equation} $I,X, Y, Z$.
For an arbitrary element $g \in G_n$,  let $g^1,...,g^n \in \scr{X}$ be such that
\begin{equation}\label{eq:canonical_tensor}
g = \alpha g^1 \otimes \hdots \otimes g^n
\end{equation}
then the sequence $g_i^1,...,g_i^n$ is the unique such, and we denote a length $2n$ sequence $x = (x_1,...,x_{2n})$ in $\bb{Z}_2^{2n}$ by $r(g_i)$ defined by the following schemata:
\begin{itemize}
\item $g^j = X$ if and only if $x_j = 1$,
\item $g^j = Z$ if and only if$x_{j + n} = 1$,
\item $g^j = Y$ if and only if $x_i = x_{j+n} = 1$.
\end{itemize}
Given a set $\lbrace g_1,...,g_k\rbrace$ of elements of the Pauli group, the \textbf{check matrix} is the $k \times 2n$ matrix whose $i^\text{th}$ row is $r(g_i)$. The check matrix is denoted $\operatorname{Check}(g_1,...,g_n)$.
\end{defn}
\begin{observation}\label{obs:evenness}
Let $(g_1,g_2)$ be a pair of elements of $G_n$ and let $g_1^1,...,g_1^n, g_2^1,...,g_2^n \in \scr{X}$ be such that
\begin{equation}
g_1 = \alpha_1 g_1^1 \otimes \hdots \otimes g_1^n,\qquad g_2 = \alpha_2 g_2^1 \otimes \hdots \otimes g_2^n
\end{equation}
we see that $g_1$ and $g_2$ commute if and only if the number of times $g_1^j$ and $g_2^j$ are distinct matrices with neither equal to the identity is even. 
\end{observation}
Defining
\begin{equation}
\Lambda :=
\begin{pmatrix}
0 & I\\
I & 0
\end{pmatrix}
\end{equation}
we have the following:
\begin{lemma}\footnote{The product is computed over $\mathbb{F}_2$ (modulo 2). See Section~\ref{sec:errata}.}
Let $(g_1,g_2) \in G_n$. Then $g_1,g_2$ commute if and only if
\begin{equation}
r(g_1) \Lambda r(g_2)^T = 0
\end{equation}
\end{lemma}
\begin{proof}[Rough sketch]
The form of $r(g_1)$:
\begin{equation}
r(g_1) =
\begin{pmatrix}
X \text{ or } Y \text{ in }g_1 & | & Z \text{ or } Y\text{ in }g_1
\end{pmatrix}
\end{equation}
and similarly for $r(g_2)$. Thus we have
\begin{equation}
r(g_1) \Lambda r(g_2)^T = \begin{pmatrix}
X \text{ or } Y \text{ in }g_1 & | & Z \text{ or } Y\text{ in }g_1
\end{pmatrix}
\begin{pmatrix}
Z \text{ or } Y \text{ in }g_2\\
X \text{ or } Y \text{ in }g_2
\end{pmatrix}
\end{equation}
This is a mathematical incarnation of the requirements specified by Observation \ref{obs:evenness}.
\end{proof}
The Check matrix is useful for more:
\begin{defn}\footnote{Index typo: the right-hand side should have $g_r$, not $g_n$. See Section~\ref{sec:errata}.}
A set of elements $g_1,...,g_r \in G_n$ of the Pauli group $G_n$ are \textbf{independent} if the for any $j$ we have, where we write $\hat{g}_i$ for the omission of $g_i$:
\begin{equation}
\langle g_1,...,g_r \rangle \neq \langle g_1,...,\hat{g}_j, ..., g_n\rangle
\end{equation}
(here, the notation $\langle g_1,...,g_n\rangle$ denotes the group generated by these elements).
\end{defn}
\begin{lemma}\label{lem:lin_indep}
Let $g_1,...,g_r \in G_n$ be a set of elements such that $-I \not\in \langle g_1,...,g_r\rangle$, then the elements $g_1,...,g_r$ are independent if and only if $r(g_1),...,r(g_r)$ and linearly independent (over the field $\bb{Z}_2$).
\end{lemma}
\begin{proof}
Exercise.
\end{proof}
Denote qubit space $\bb{C}^2$ by $\bb{H}$. Given a subgroup $S \subseteq G_n$ of the Pauli group $G_n$, we denote by $V_S$ the set of vectors $\ket{\psi} \in \bb{H}^{\otimes n}$ such satisfying
\begin{equation}
\forall W \in S, W\ket{\psi} = \ket{\psi}
\end{equation}
The following Lemma will be used to calculate the dimension of $V_S$:
\begin{lemma}\label{lem:anti_commutes}\footnote{The proof says ``$k$ linearly independent columns''; should read ``rows''. The vector $x$ should lie in $\mathbb{F}_2^{2n}$, not $\mathbb{F}_2^k$ (the wrong dimension). See Section~\ref{sec:errata}.}
Let $g_1,...,g_k$ be independent elements of the Pauli group $G_n$ and denote by $S$ the group they generate. Assume $-I \not\in S$. Then for each $i = 1,...,k$ there exists $g \in G_n$ such that $g$ anti-commutes with $g_i$ and commutes with all $g_j$ such that $i \neq j$.
\end{lemma}
\begin{proof}
The set $r(g_1),...,r(g_k)$ is linearly independent by Lemma \ref{lem:lin_indep}, thus the check matrix of $g_1,...,g_k$ has $k$ linearly independent columns. So, there exists a vector $x \in \bb{Z}_2^k$ such that
\begin{equation}
\operatorname{Check}(g_1,...,g_n)\Lambda x = e_i
\end{equation}
where $e_i$ is the $i^{\text{th}}$ standard basis vector of $\bb{Z}_2^k$. Let $g$ be such that $r(g)^T = x$. The result follows from Lemma \ref{lem:lin_indep}.
\end{proof}
\begin{thm}\footnote{The conclusion is correct, but the proof has issues: ``$I$ is a projector onto an $n$-dimensional space'' should read ``$2^n$-dimensional'' (since $n$ qubits give $(\mathbb{C}^2)^{\otimes n}$ of dimension $2^n$). Independence requires $g_j$ Hermitian with eigenvalues $\pm 1$, which uses $-I\notin S$ to rule out phase-$i$ generators. See Section~\ref{sec:errata}.}
Let $S = \langle g_1,...,g_k\rangle$ and say $-I \not\in S$. Then $\operatorname{dim}V_S = 2^{n - k}$.
\end{thm}
\begin{proof}
We notice that $(1/2)(I + g_j)$ is the projector onto the +1-Eigenspace of $g_j$. We let $x = (x_1,...,x_k) \in \bb{Z}_2^k$ and define the operator
\begin{equation}
P_S^x := 1/2^{k}\prod_{j = 1}^k(I + (-1)^{x_{j}}g_j)
\end{equation}
By Lemma \ref{lem:anti_commutes} we have for each $g_j$ there exists $g_{x_j}$ such that $g_{x_j}g_j g_{x_j}^{-1} = -g_j$. Let $g_x = g_{x_1}\hdots g_{x_k}$, then
\begin{align*}
g_xP_S^{(0,...,0)} g_x^{-1} &= 1/2^k\prod_{j = 1}^k (g_{x_j}g_{x_j}^{-1} + g_{x_j} g_{j}g_{x_j}^{-1})\\
&= P_S^x
\end{align*}
Thus there is an isomorphism
\begin{equation}
\operatorname{im}P_S^x \cong \operatorname{im}P_S^{(0,...,0)}
\end{equation}
Since $\operatorname{im}P_S \cong V_S$ we have $\operatorname{dim}\operatorname{im}P_S^x = \operatorname{dim}V_S$. Finally we note that
\begin{equation}
I = \sum_{x \in \bb{Z}_2^k}P_S^x
\end{equation}
The operator $I$ is a projector onto an $n$-dimensional space, and $\sum_{x \in \bb{Z}_2^k}P_S^x$ is a sum of $2^k$ orthogonal projectors all of the same dimension as $V_S$, thus
the only possibility is $\operatorname{dim}V_S = 2^{n - k}$.
\end{proof}

\clearpage
\section{AI-generated errata}\label{sec:errata}

This section was produced by a large language model. It corrects statements and proof steps in the body. Each subsection mirrors a section above; corrected statements supersede the originals.

\subsection{Mathematical definitions}

\subsubsection*{Qubit}

The standard convention is: $\mathbb{C}^2$ is a single-qubit state space, and $(\mathbb{C}^2)^{\otimes n}$ is the $n$-qubit state space. The opening definition should read ``an $n$-qubit state space is $(\mathbb{C}^2)^{\otimes n}$''.

\subsubsection*{Trace}

For the formula $\operatorname{Tr}(A)=\sum_i\bra{e_i}A\ket{e_i}$ to hold, the basis $\{\ket{e_i}\}$ must be \emph{orthonormal}, not merely orthogonal. For an orthogonal basis $\{\ket{v_i}\}$ the formula carries norm factors. Throughout the document, finite-dimensionality of $\mathbb{H}$ should be assumed at the start, since the trace is otherwise only defined for trace-class operators.

\subsubsection*{Measurements}

The completeness condition is $\sum_{m\in\mathcal{M}}M_m^\dagger M_m=I$, with subscripts on each $M_m$. The post-measurement state $M_m\ket{\psi}/\sqrt{p(m)}$ is defined only when $p(m)>0$.

\subsubsection*{Projectors}

A projector in quantum mechanics is an \emph{orthogonal} projection: $P^2=P$ and $P^\dagger=P$. The bare condition $P^2=P$ defines an algebraic idempotent, which is not what projective measurements use. The projectors $P_m$ from the spectral theorem of a Hermitian operator are automatically orthogonal.

\subsubsection*{Extension to a unitary}

The lemma stating that $U:W\to V$ extends to a unitary is incorrectly stated. If $U$ is literally unitary onto $V$ then $\dim W=\dim V$, forcing $W=V$. The intended hypothesis is that $U$ is an \emph{isometry} $W\to V$.

\begin{lemma}
Fix a finite-dimensional Hilbert space $V$, a subspace $W\subseteq V$, and an isometry $U:W\to V$. Then $U$ extends to a unitary $U':V\to V$.
\end{lemma}
\begin{proof}
$U(W)\subseteq V$ has dimension $\dim W$. Choose any unitary $W^\perp\to U(W)^\perp$, and define $U'$ to equal $U$ on $W$ and the chosen unitary on $W^\perp$. Then $U'$ is unitary on $V$.
\end{proof}

The original proof $U'=\operatorname{Proj}_W\circ U$ is wrong: domain and codomain do not match, and the result is not a unitary extension.

\subsubsection*{Naimark dilation}

The proposition is essentially correct but the proof has many notation and tensor-order errors. A clean version:

\begin{proposition}
Fix a measurement $\{M_m\}_{m\in\mathcal{M}}$ on $\mathbb{H}$. Let $Q$ have orthonormal basis $\{\ket{m}:m\in\mathcal{M}\}$ and define $V:\mathbb{H}\to\mathbb{H}\otimes Q$ by
\[
V\ket{\psi}=\sum_m M_m\ket{\psi}\otimes\ket{m}.
\]
Then $V^\dagger V=\sum_m M_m^\dagger M_m=I$, so $V$ is an isometry. Fix $\ket{\varphi}\in Q$ and extend
\[
U(\ket{\psi}\otimes\ket{\varphi})=\sum_m M_m\ket{\psi}\otimes\ket{m}
\]
to a unitary $U$ on $\mathbb{H}\otimes Q$ (using the corrected isometry-extension lemma). Measuring the ancilla with projectors $P_n=I_\mathbb{H}\otimes\ket{n}\bra{n}$ on $U(\ket{\psi}\otimes\ket{\varphi})$ gives
\[
p(n)=\bra{\psi}M_n^\dagger M_n\ket{\psi}.
\]
\end{proposition}

The original proof switches between $Q\otimes\mathbb{H}$ and $\mathbb{H}\otimes Q$, writes $I_q\otimes\ket{m}\bra{m}$ where $I_\mathbb{H}\otimes\ket{m}\bra{m}$ is intended, omits subscripts on $M_m^\dagger M_{m'}$, and has a wrong final ket in the inner-product calculation.

\subsection{Density operator}

\subsubsection*{Pure versus mixed}

Purity is a property of the density operator, not the ensemble. The density operator $\rho$ is \emph{pure} iff $\rho=\ket{\psi}\bra{\psi}$ for some unit vector, equivalently $\operatorname{Tr}(\rho^2)=1$. An ensemble of length $>1$ can still give a pure $\rho$.

\subsubsection*{Measurement update (operator order)}

The post-measurement density operator is
\[
\rho_m=\frac{M_m\rho M_m^\dagger}{\operatorname{Tr}(M_m^\dagger M_m\rho)},
\]
not $M_m^\dagger\rho M_m$. With the correct ordering, the pure-state and density-operator measurement rules agree:
\[
\frac{M_m\ket{\psi}\bra{\psi}M_m^\dagger}{p(m)}=\Big(\frac{M_m\ket{\psi}}{\sqrt{p(m)}}\Big)\Big(\frac{M_m\ket{\psi}}{\sqrt{p(m)}}\Big)^\dagger.
\]

The displayed expression $\rho_m=\sum_{i=1}^n \frac{M_m\rho M_m^\dagger}{\operatorname{Tr}(M_m^\dagger M\rho)}$ in the derivation also has an erroneous factor of $n$; the sum disappears once $\rho$ is substituted.

\subsubsection*{Time evolution}

The first evolved state in the sequence is $U_1\rho U_1^\dagger$, not $U\rho U^\dagger$:
\[
(\rho,\ U_1\rho U_1^\dagger,\ U_2 U_1\rho U_1^\dagger U_2^\dagger,\ \ldots,\ U_n\cdots U_1\rho U_1^\dagger\cdots U_n^\dagger).
\]

\subsubsection*{Equivalence of ensembles}

The proposition as written is mathematically incoherent: the notation $\{p_i\ket{i}\bra{i}\}$ is not standard ensemble notation, the indices and bras are conflated, and the proof omits bras in spectral decompositions. The correct theorem is the standard unitary-freedom-of-ensembles theorem.

\begin{proposition}
Fix two ensembles $\{p_i,\ket{\psi_i}\}_{i=1}^m$ and $\{q_j,\ket{\varphi_j}\}_{j=1}^r$. Define unnormalised vectors
\[
\ket{\widetilde{\psi}_i}=\sqrt{p_i}\ket{\psi_i},\qquad \ket{\widetilde{\varphi}_j}=\sqrt{q_j}\ket{\varphi_j},
\]
and pad the shorter list with zero vectors so both have the same length $N$. The two ensembles determine the same density operator iff there exists a unitary $U=(u_{ij})$ with
\[
\ket{\widetilde{\psi}_i}=\sum_j u_{ij}\ket{\widetilde{\varphi}_j}.
\]
\end{proposition}

For ensembles of unequal length without padding, the appropriate object is an isometry or partial isometry rather than a square unitary.

The preceding observation also has errors. The diagonal form is
\[
\sum_i p_i\ket{\psi_i}\bra{\psi_i}=\sum_j\lambda_j\ket{\varphi_j}\bra{\varphi_j},
\]
not $\sum_j\lambda_j\ket{\varphi_j}$ (a missing bra). Coefficient comparison gives
\[
\lambda_j\delta_{jk}=\sum_i p_i\alpha_{ij}\overline{\alpha_{ik}},
\]
not $\sum_i p_i\alpha_{ij}\alpha_{ij}^\dagger$.

\subsection{Partial trace}

The intermediate calculation in the Bell-state example duplicates the $\ket{11}\bra{00}$ line. The correct partial traces are:
\[
\operatorname{Tr}_2(\ket{11}\bra{00})=\operatorname{Tr}(F_{01})E_{01}=0,\quad
\operatorname{Tr}_2(\ket{00}\bra{11})=\operatorname{Tr}(F_{10})E_{10}=0,\quad
\operatorname{Tr}_2(\ket{11}\bra{11})=\operatorname{Tr}(F_{11})E_{11}=\ket{1}\bra{1}.
\]
The final result $\operatorname{Tr}_2(\rho)=I/2$ is correct.

\subsubsection*{Adjunction}

The standard finite-dimensional tensor-Hom adjunction is
\[
\operatorname{Hom}(V\otimes U,W)\cong\operatorname{Hom}(U,V^*\otimes W),
\]
not $\operatorname{Hom}(V,W\otimes V^*)$. The right adjoint of $V\otimes-$ is $V^*\otimes-$.

\subsection{Quantum operations}

A quantum operation is a linear map $\mathcal{E}:L(\mathbb{H}_p)\to L(\mathbb{H}_p)$ on operators, not $\mathbb{H}_p\to\mathbb{H}_p$. Trace-preserving completely positive such maps are called quantum channels.

\subsubsection*{Kraus form}

For a pure environment state $\rho_e=\ket{e_0}\bra{e_0}$, the Kraus operators are
\[
E_i=(I\otimes\bra{i})U(I\otimes\ket{e_0}),
\]
giving $\mathcal{E}(\rho)=\sum_i E_i\rho E_i^\dagger$. The original derivation conflates the summation index with the fixed environment state, writes $\rho\otimes\ket{i}\bra{i}$ where $\rho\otimes\ket{e_0}\bra{e_0}$ is intended, and uses an unnecessary multiplication map.

\begin{lemma}
A trace-preserving quantum operation $\mathcal{E}$ admits an operator-sum representation $\mathcal{E}(\rho)=\sum_i E_i\rho E_i^\dagger$ with $\sum_i E_i^\dagger E_i=I$. The Kraus operators $\{E_i\}$ are not unique.
\end{lemma}

For mixed environment $\rho_e=\sum_a p_a\ket{e_a}\bra{e_a}$, the Kraus operators are $E_{ia}=\sqrt{p_a}(I\otimes\bra{i})U(I\otimes\ket{e_a})$.

\subsection{Quantum error correction}

\subsubsection*{Basis convention}

The standard convention (consistent with $Z=\operatorname{diag}(1,-1)$) is $\ket{0}=(1,0)^T$ and $\ket{1}=(0,1)^T$, giving $Z\ket{0}=\ket{0}$ and $Z\ket{1}=-\ket{1}$. The body's reversed convention contradicts later calculations and should be corrected.

\subsubsection*{Bell states versus Hadamard-basis states}

The states $\ket{\pm}=(\ket{0}\pm\ket{1})/\sqrt{2}$ are the Hadamard-basis (or $X$-basis) states, not Bell states. Bell states are two-qubit entangled states such as $(\ket{00}+\ket{11})/\sqrt{2}$.

\subsubsection*{Encoding}

A quantum encoding $\iota:\mathbb{H}\to\mathbb{H}^{\otimes n}$ should be a linear isometry, not merely an injective map, in order to preserve superpositions and inner products.

\subsubsection*{Bit-flip correction: measurement versus application}

The proof confuses two distinct operations. Writing $\bra{\psi}Z_1Z_2\ket{\psi}$ gives the expectation value of the observable, not the projective measurement. The projective measurement of $Z_1Z_2$ uses projectors
\[
P_\pm=\frac{I\pm Z_1Z_2}{2},
\]
and the post-measurement state is $P_\pm\ket{\psi}/\sqrt{p(\pm)}$. Code states are eigenvectors of $Z_1Z_2,Z_2Z_3$, so the syndrome measurement reveals the syndrome without disturbing the state. Applying $Z_1Z_2$ as a unitary is a different operation entirely.

The amplitude error $-a^2-b^2=-1$ should read $-|a|^2-|b|^2=-1$ for general $a,b\in\mathbb{C}$. The bra should be $\overline{a}\bra{010}+\overline{b}\bra{101}$.

\subsubsection*{Correctable set of errors}

For the bit-flip code, the correctable errors are $\{I,X_1,X_2,X_3\}$, not $\{Z_1Z_2,Z_1Z_3,Z_2Z_3\}$. The latter are syndrome observables. A correct recovery uses syndrome projectors $P_{\pm\pm}$ and corrections $I,X_1,X_3,X_2$ for syndromes $(++,-+,+-,--)$ respectively.

\subsubsection*{Phase-flip}

The same measurement-versus-application confusion appears. The projective measurements use $(I\pm X_1X_2)/2$ and $(I\pm X_2X_3)/2$. The identity $H^{\otimes n}Z_iH^{\otimes n}=X_i$ is correct.

\subsection{Shor code}

\subsubsection*{Encoding type}

The composition $\operatorname{BitFlip}\circ\operatorname{PhaseFlip}$ is not type-correct: $\operatorname{PhaseFlip}:\mathbb{H}\to\mathbb{H}^{\otimes 3}$ produces three qubits, but $\operatorname{BitFlip}:\mathbb{H}\to\mathbb{H}^{\otimes 3}$ accepts a single qubit. The correct composition is
\[
\operatorname{Shor}=\operatorname{BitFlip}^{\otimes 3}\circ\operatorname{PhaseFlip}.
\]
Explicitly,
\[
\ket{0}\mapsto\frac{1}{2\sqrt{2}}(\ket{000}+\ket{111})^{\otimes 3},\qquad \ket{1}\mapsto\frac{1}{2\sqrt{2}}(\ket{000}-\ket{111})^{\otimes 3}.
\]

\subsubsection*{Stabiliser measurements}

The Shor code does not use $Z_iZ_{i+1}$ for $i=1,\ldots,8$. The standard $Z$-type stabiliser generators are
\[
Z_1Z_2,\ Z_2Z_3,\ Z_4Z_5,\ Z_5Z_6,\ Z_7Z_8,\ Z_8Z_9,
\]
checking bit flips within each block of three. Cross-block checks $Z_3Z_4,Z_6Z_7$ are not stabilisers and disturb the encoded information.

The standard $X$-type stabiliser generators are
\[
X_1X_2X_3X_4X_5X_6,\ X_4X_5X_6X_7X_8X_9.
\]
This gives 8 stabiliser generators total, not 10.

\subsubsection*{Recovery table}

The original recovery table has length-9 rows where length-8 was promised, repeats rows, has impossible syndrome patterns, and mixes bit-flip and phase-flip recovery inconsistently. The correct rule:
\begin{itemize}
\item Use the six $Z$-type checks to identify the bit-flip location and apply $X_i$.
\item Use the two $X$-type checks to identify the block with a phase flip; apply $Z$ to any qubit in that block (or equivalently $Z_iZ_{i+1}Z_{i+2}$ to the whole block).
\end{itemize}
The phase-syndrome interpretation:
\[
\begin{array}{c|c}
(s_1,s_2) & \text{phase-error block}\\\hline
(+1,+1) & \text{none}\\
(-1,+1) & \text{block 1 (qubits 1,2,3)}\\
(-1,-1) & \text{block 2 (qubits 4,5,6)}\\
(+1,-1) & \text{block 3 (qubits 7,8,9)}
\end{array}
\]

\subsubsection*{Combined errors}

For a combined error $Z_iX_j\ket{\varphi}$, the correction is correct up to a global phase. This requires using that Pauli operators either commute or anticommute, and anticommutation contributes only an unobservable global phase after both corrections.

\subsection{Stabiliser formalism}

\subsubsection*{Pauli group}

The cleanest statement:
\[
G_n=\{\alpha P_1\otimes\cdots\otimes P_n:\alpha\in\{\pm 1,\pm i\},\ P_j\in\{I,X,Y,Z\}\}.
\]

\subsubsection*{Binary representation}

For $g=\alpha g^1\otimes\cdots\otimes g^n$, the binary vector $r(g)=(x_1,\ldots,x_n,z_1,\ldots,z_n)$ satisfies
\[
x_j=1\iff g^j\in\{X,Y\},\qquad z_j=1\iff g^j\in\{Y,Z\}.
\]
The original ``$g^j=X$ iff $x_j=1$'' fails for $g^j=Y$, since $Y$ also has $x_j=1$. The single-qubit map is $I\mapsto(0,0)$, $X\mapsto(1,0)$, $Z\mapsto(0,1)$, $Y\mapsto(1,1)$.

There is also an index typo: ``$x_i=x_{j+n}=1$'' should read ``$x_j=x_{j+n}=1$''.

\subsubsection*{Symplectic commutation}

The condition $r(g_1)\Lambda r(g_2)^T=0$ should be computed over $\mathbb{F}_2$ (i.e.\ modulo 2), with
\[
\Lambda=\begin{pmatrix}0 & I\\ I & 0\end{pmatrix}.
\]

\subsubsection*{Independence}

The notation in the independence definition has a typo: $\langle g_1,\ldots,\widehat{g}_j,\ldots,g_n\rangle$ should be $\langle g_1,\ldots,\widehat{g}_j,\ldots,g_r\rangle$. The hypothesis $-I\notin\langle g_1,\ldots,g_r\rangle$ is essential for binary-vector independence to coincide with Pauli-group independence.

\subsubsection*{Anticommuting Pauli}

The check matrix $H$ is $k\times 2n$. Linear independence of the $r(g_i)$ means $H$ has $k$ linearly independent \emph{rows} (not columns), so $\operatorname{rank}H=k$. Since $\Lambda$ is invertible, $H\Lambda$ has rank $k$, so $x\mapsto H\Lambda x$ from $\mathbb{F}_2^{2n}$ to $\mathbb{F}_2^k$ is surjective. For each $e_i\in\mathbb{F}_2^k$ there exists $x\in\mathbb{F}_2^{2n}$ with $H\Lambda x=e_i$. The original ``$x\in\mathbb{Z}_2^k$'' has the wrong dimension.

\subsubsection*{Stabiliser dimension theorem}

The conclusion $\dim V_S=2^{n-k}$ is correct. The proof needs repair.

Each $g_j$ is Hermitian with eigenvalues $\pm 1$ (since $-I\notin S$, no element has order-4 phase). The simultaneous eigenspace projectors are
\[
P_x=\frac{1}{2^k}\prod_{j=1}^k\big(I+(-1)^{x_j}g_j\big).
\]
For each $j$, choose $h_j\in G_n$ anticommuting with $g_j$ and commuting with $g_\ell$ for $\ell\neq j$ (Lemma~\ref{lem:anti_commutes}); for $x=(x_1,\ldots,x_k)$ set $h_x=\prod_{j:x_j=1}h_j$. Then conjugation by $h_x$ sends $P_0$ to $P_x$, so all $P_x$ have equal rank.

The physical Hilbert space is $(\mathbb{C}^2)^{\otimes n}$, dimension $2^n$ (the original ``the operator $I$ is a projector onto an $n$-dimensional space'' is wrong: $n$ qubits give a $2^n$-dimensional space). Since $I=\sum_x P_x$ is a sum of $2^k$ orthogonal projectors of equal rank summing to $2^n$,
\[
\dim V_S=\operatorname{rank}P_0=\frac{2^n}{2^k}=2^{n-k}.
\]

\begin{thebibliography}{9}
\bibitem{quantum_computing} M. Nielsen and I. Chuang, \emph{Quantum Computation and Quantum Information}, Cambridge University Press, 2010.

\bibitem{analysis} W. Troiani \emph{Analysis}

\bibitem{CommutativeAlgebra} W. Troiani \emph{Commutative Algebra}
\end{thebibliography}
\end{document}




































\maketitle
\tableofcontents
